\documentclass[11pt,a4paper]{article}
\usepackage[margin=24mm]{geometry}
\usepackage{amsmath,amssymb,amsthm,array}
\setlength{\emergencystretch}{2em}
\newtheorem{theorem}{Theorem}
\newcommand{\gen}[1]{\langle#1\rangle}
\title{The third prefix at a singleton end}
\author{Human Compression project: L044 continuation}
\date{}
\begin{document}\maketitle
Retain the critical singleton hypotheses, the actual natural omission
$G_0=\operatorname{Alt}(\Omega\setminus\{a\})$, and the complete
prefix-growth proof. Thus all marks after $\rho_0$ fix $a$,
$C_k=a^{G_{\le k}}$ grows strictly, and distant later marks fix
$C_k$ pointwise. Actual ambient IP is assumed throughout.

\begin{theorem}
If $P_2=\gen{\rho_0,\rho_1,\rho_2}$ is intransitive on its
support, that support has at least eight points.
Consequently a first shared label $g=3$ cannot have an
intransitive prefix in the critical singleton support interface.
\end{theorem}
\begin{proof}
Suppose $P_2$ moves at most seven points. Pad this set with
fixed points if necessary, using $n\ge14$, and name $a=1$.
Every nonidentity even involution on seven points is a double
transposition. Since $\rho_0$ moves $a$, we may normalize
\[
                    \rho_0=(1\ 2)(3\ 4).
\]
The mark $\rho_2$ fixes $1$ and commutes with $\rho_0$,
so fixes $2$ and preserves $\{3,4\}$. On the other three
points an involution has at most one transposition. Evenness
and nonidentity therefore force, after renaming these points,
\[
                    \rho_2=(3\ 4)(5\ 6),
\]
with $7$ fixed. Strict growth $C_0\subsetneq C_1$ makes
$\rho_1$ move $2$, while it fixes $1$. Write
\[
                    \rho_1=(2\ x)(y\ z).
\]
The independent interchanges of $3,4$ and of $5,6$ preserve
both displayed marks, so we may take $x\in\{3,5,7\}$.
The pair $\{y,z\}$ is any two-element subset of the remaining
four points of $\{2,\ldots,7\}$. These are exactly the
eighteen cases below; no group census is needed.

\begin{center}
\begin{tabular}{c|l|l}
$x$ & possible $\{y,z\}$ & consequence\\\hline
$3$ & $\{4,5\},\{4,6\}$ & $P_2$ transitive on its support\\
$3$ & $\{4,7\},\{5,6\},\{5,7\},\{6,7\}$ & $C_1=C_2$\\
$5$ & $\{3,6\},\{4,6\}$ & $P_2$ transitive on its support\\
$5$ & $\{3,4\}$ & preserved exterior pair\\
$5$ & $\{3,7\},\{4,7\}$ & preserved exterior triple\\
$5$ & $\{6,7\}$ & actual point-stabilizer contradiction\\
$7$ & all six possible pairs & $C_1=C_2$
\end{tabular}
\end{center}
Here are the details of every row. In the first and third rows,
the transpositions of the three marks join precisely
$\{1,\ldots,6\}$, and $7$ is fixed. Their action is therefore
transitive on the entire support, contrary to the hypothesis.
In the second row, $C_1=C_2$ is $\{1,2,3,4,7\}$ for the
pair $\{4,7\}$ and is $\{1,2,3,4\}$ for the other three
pairs. In the last row, $\rho_2$ fixes $C_1=\{1,2,7\}$
pointwise, again giving $C_1=C_2$. Strict growth excludes
all these cases.

For $x=5$, $\{y,z\}=\{3,4\}$, the two nontrivial
$P_2$-orbits are
\[
                  C_2=\{1,2,5,6\},\qquad E=\{3,4\}.
\]
The $P_1$-orbits are $C_1=\{1,2,5\}$, $E$, and fixed
singletons. Thus $\rho_3$, which centralizes $P_1$, preserves
the unique two-point orbit $E$. Every mark of index at least
four centralizes $P_2$ and likewise preserves $E$.
The first three marks preserve it by its definition as a
$P_2$-orbit. Hence $E$ is a proper $G$-invariant set,
contrary to transitivity.

For $x=5$ and $\{y,z\}=\{3,7\}$ or $\{4,7\}$,
the two nontrivial $P_2$-orbits are
\[
                  C_2=\{1,2,5,6\},\qquad E=\{3,4,7\}.
\]
The two nontrivial $P_1$-orbits are $C_1=\{1,2,5\}$ and
$E$. The mark $\rho_3$ centralizes $P_1$ and fixes $1$,
so fixes $C_1$ pointwise. It therefore cannot exchange these
two three-point orbits, and must preserve $E$.
All marks of index at least four centralize $P_2$, whose
nontrivial orbits have distinct sizes four and three, so also
preserve $E$. Once again $G$ preserves $E$, a contradiction.

Finally suppose $\rho_1=(2\ 5)(6\ 7)$. On
$\{1,2,5\}$ the product $\rho_0\rho_1$ is a three-cycle,
and off this set it is $(3\ 4)(6\ 7)$. Thus
\[
             w=(\rho_0\rho_1)^3=(3\ 4)(6\ 7).
\]
This is a nonidentity even permutation fixing $a=1$,
so belongs to the actual natural group $G_0$. It also
belongs to $P_1$. Full IP gives
\[
               w\in P_1\cap G_0=\gen{\rho_1}.
\]
But $w$ fixes $2$, whereas $\rho_1$ moves it. This is
impossible. All eighteen choices are exhausted.

The first-shared support theorem gives
$|\operatorname{supp}(P_{g-1})|\le2g+1$.
At $g=3$ its intransitive-prefix alternative would give
the forbidden support bound seven, proving the consequence.
\end{proof}

\paragraph{Attribution and boundary.}
The complete singleton-prefix growth and support proofs are retained.
The finite normalization here is a direct hand argument on three
literal permutations; an exploratory enumeration suggested its case
split but is not a proof premise. The point-stabilizer step uses
the acquired actual natural omission and full IP, and does not
assign IP or fidelity to any orbit projection. The source route is
the singleton continuation of Cameron--Fernandes--Leemans--Mixer,
\emph{Highest rank of a polytope for $A_n$}, Propositions~5.28--5.32.
Larger first shared indices and complete singleton routing remain.
No whole-original, question, numerical or formal credit is claimed.
\end{document}
